\section{Method}
\label{sec:method}

\subsection{Compass plane}
\label{sec:method:plane}

For attention head $h$ at layer $\ell$ we form
$W_{OV}^{(\ell,h)} = W_V^{(\ell,h)} W_O^{(\ell,h)}$ and compute its
SVD $W_{OV} = U \Sigma V^\top$, indexed by $\sigma_1 \geq \sigma_2
\geq \cdots \geq 0$. Following the convention in
\S\ref{sec:background} we use rows of $V^\top$ as write-out
directions and write $u_i := V^\top_{[i,:]}$ throughout. A
\emph{compass plane} is an ordered pair
$\mathcal{C}_{i,j} = \mathrm{span}(u_i, u_j)$ with $i \neq j$.

\subsection{Causal $\alpha$-sweep}
\label{sec:method:sweep}

The compass test is a causal injection that varies an angle
$\theta$ and a scale $\alpha$. \cref{alg:sweep} states the
procedure; \cref{eq:alpha_sweep} gives the injection.

\begin{equation}
v(\theta, \alpha) =
  \alpha\,\sigma_i\,\cos\theta\,u_i +
  \alpha\,\sigma_j\,\sin\theta\,u_j.
\label{eq:alpha_sweep}
\end{equation}

\begin{algorithm}[h]
\small
\caption{Causal $\alpha$-sweep on a compass plane.}
\label{alg:sweep}
\begin{algorithmic}[1]
\Require model $M$, layer $\ell$, head $h$, dim pair $(i, j)$,
probe tokens $(t_+, t_-)$, prompt set $\mathcal{P}$,
angles $\Theta = \{0, \tfrac{2\pi}{36}, \dots\}$,
scales $\mathcal{A} = \{1, 3, 10\}$
\State $W_{OV} \gets W_V^{(\ell,h)} W_O^{(\ell,h)}$;
$\;U \Sigma V^\top \gets \mathrm{SVD}(W_{OV})$
\State $u_i, u_j \gets V^\top_{[i,:]}, V^\top_{[j,:]}$;
$\;\sigma_i, \sigma_j \gets \Sigma_{ii}, \Sigma_{jj}$
\For{$\alpha \in \mathcal{A}$, $\theta \in \Theta$, $p \in \mathcal{P}$}
  \State Hook \texttt{blocks.}$\ell$\texttt{.hook\_resid\_pre}:
         $\;r \gets r + v(\theta, \alpha)$ on the last token only
  \State Forward pass; record
         $\mathrm{LD}(\theta, \alpha; p) = \mathrm{logit}_{t_+} -
          \mathrm{logit}_{t_-}$
\EndFor
\State Average $\mathrm{LD}$ over $\mathcal{P}$
\State Fit $\mathrm{LD}(\theta) \approx \mu + A\cos(\theta - \phi)$
       via first DFT bin
\State \Return $(\mu, A, \phi)$ for each $\alpha$
\end{algorithmic}
\end{algorithm}

We average each $\mathrm{LD}$ across three prompts per condition
(neutral, $+$-context, $-$-context), so each $\alpha$-sweep produces
a 36-point curve from which $(\mu, A, \phi)$ are read.

\paragraph{The two pillars.}
A causal compass passes both: amplitude linearity ($A(\alpha)
\propto \alpha$, $R^2 \geq 0.95$ through the origin) and phase
stability ($|\phi(\alpha_{\max}) - \phi(\alpha_{\min})| \leq 10^\circ$).
Phase stability is the discriminating test: a generic 2D direction
in the residual will show some amplitude under linear injection
(the linearity pillar partly tracks the linear injection itself),
but only a direction the head actually writes through produces a
phase that is invariant to scale. We require both pillars together
because the joint constraint is rare in random subspaces: random
orthonormal pairs in our scans pass linearity routinely but rarely
hold phase below $10^\circ$.

\subsection{Falsification battery}
\label{sec:method:battery}

We complement \cref{eq:alpha_sweep} with five pre-specified tests
(verdicts in parentheses; $\checkmark$ = signal preserved,
$\times$ = signal killed as a falsification target):
\textbf{decode-heads} ($\checkmark$): $W_U^\top u_i$ promotes
readable categorical tokens at $\theta \in \{0, \pi/2, \pi,
3\pi/2\}$;
\textbf{cyclicity} ($\checkmark$): top-10 decode at $\theta = 0$
matches $\theta = 2\pi$ and is Jaccard-disjoint from $\theta = \pi$;
\textbf{permutation null} ($\times$): permuting $\pm$-context labels
across prompts must kill the signal --- this rules out the
artefactual reading that the sinusoidal response merely tracks the
prompt set rather than the categorical contrast;
\textbf{random-plane null} ($\times$): random orthonormal directions
in $\mathbb{R}^{d_{\text{model}}}$ should remove the causal effect;
\textbf{causal-patch} ($\checkmark$): rotating the residual's
\emph{natural} in-plane projection by $\Delta \in \{\pm 90^\circ,
\pm 180^\circ\}$ preserves magnitude and shifts logits in the
predicted direction.

The random-plane null is the most discriminating of the five in
principle; \S\ref{sec:results:null} reports its outcome across the
four-model panel and uses the result to scope the central claim.

\subsection{Models and target heads}
\label{sec:method:models}

We use TransformerLens \citep{nanda2022transformerlens} with
\texttt{HookedTransformer} on bfloat16. Targets:

\begin{itemize}
\itemsep0pt
\item \textbf{GPT-2 Small} \citep{radford2019language}: L9H7 $(u_1, u_2)$
(decode-discovered) and L10H9 $(u_0, u_3)$ (scan-discovered).
\item \textbf{Phi-3 Mini} \citep{abdin2024phi3}: L24H10 $(u_0, u_1)$
(storage-dominant); L28H1 $(u_0, u_3)$ (steering-dominant);
L24H17 $(u_1, u_2)$ (temporal); L24H28 $(u_3, u_7)$ (entity).
\item \textbf{Gemma-2-2B} \citep{team2024gemma}: L21H4 $(u_0, u_2)$.
\item \textbf{Llama-3.2-3B} \citep{grattafiori2024llama3}:
L26H14 $(u_2, u_3)$ (scan-discovered).
\end{itemize}

The blind-scan amplitude threshold is $A_{\text{slope}} \geq 0.20$
for GPT-2, Phi-3, and Llama, and $0.08$ for Gemma. Gemma's residual
is normalized more aggressively after pre-norm, so per-unit $\alpha$
writes a smaller fraction of the residual; using the model-specific
threshold recovers the same head as the unit-norm injection. The
single-head story does not change at $A_{\text{slope}} = 0.20$
either; the lower threshold widens the scan, it does not flip the
identity of the discovered head.
